\subsection{The full low-genus completion range and two source corrections}

\begin{result}{D.19}{Known; transfer proved}{Congruence comparison in genus at most two}
For every stable type of genus $0,1,$ or $2$, the natural map
$\widehat\Gamma_{g,n}\to\check\Gamma_{g,n}$
is an isomorphism. Consequently the full finite-level and procongruence completions of every $P_k$ are canonically isomorphic in these genera. In particular, their $P_1$ are connected and their $P_k$, $k\geq2$, have trivial profinite fundamental group. The full unmodified curve complex is connected with trivial profinite fundamental group in this range exactly when $d\geq3$; the two connected complexity-two exceptions have nontrivial profinite fundamental group by D.12 and D.21.
\end{result}
\begin{proof}
The genus-zero congruence theorem is the puncture induction in Boggi's Proposition~2.6, reproducing Asada's argument: start with $\Gamma_{0,3}=1$, use the completed Birman sequence, and use the tangential lift to identify the point-pushing closure with the full profinite surface group \cite{D_Boggi2009}. The opened Theorem~3.5 and Lemmas~3.6--3.7 of the same paper prove faithfulness for the full hyperelliptic group with at least one puncture. In genera one and two the hyperelliptic locus is the entire moduli stack, so this gives every stable genus-one type and every genus-two type with $n\geq1$. The closed genus-two case requires an additional argument; it is not part of that open-surface theorem.

Here is the independent closed-surface transfer, which avoids both problematic claims discussed in D.20. Write
\[
 P=\widehat{\pi_1(S_2)},\qquad
 Q=\widehat{\pi_1(\operatorname{Conf}_2(S_2))},\qquad
 F=\widehat{\pi_1(S_2\setminus\{x\})}.
\]
The full universal-curve and configuration sequences give
\[
 1\longrightarrow Q\longrightarrow\widehat\Gamma_{2,2}
       \longrightarrow\widehat\Gamma_{2,0}\longrightarrow1,
 \qquad
 1\longrightarrow F\longrightarrow Q\longrightarrow P\longrightarrow1.
\]
These are sequences of full geometric fundamental groups, obtained from the ordered configuration family, as in Part~E; no congruence exactness is assumed. Its geometric monodromy preserves the forgetful fibrations and diagonal inertia, so its image in $\operatorname{Out}(Q)$ is FC-admissible in the sense of Hoshi--Mochizuki. The injectivity assertion of their Theorem~6.1, with all primes and $n=1$, applies to this closed genus-two curve and gives
\[
 \operatorname{Out}^{\rm FC}(Q)\hookrightarrow
 \operatorname{Out}^{\rm FC}(P).
\]
We opened that theorem and its proof, including its reduction to Corollary~5.3; it allows a proper hyperbolic curve at this value of $n$ \cite{D_HoshiMochizuki2011}.

Let $y\in\widehat\Gamma_{2,0}$ act trivially in $\operatorname{Out}(P)$. The displayed injection makes its outer action on $Q$ trivial. Lift $y$ to $z\in\widehat\Gamma_{2,2}$. Conjugation by $z$ on $Q$ is inner, so multiplication by a suitable element of $Q$ makes $z$ centralize $Q$, without changing its image $y$. In particular $z$ centralizes the distinguished subgroup $F$. Its image $z_1\in\widehat\Gamma_{2,1}$ therefore acts trivially in $\operatorname{Out}(F)$: this action is precisely restriction of conjugation by a lift in $\widehat\Gamma_{2,2}$. The already established open genus-two congruence theorem forces $z_1=1$, and its image $y$ is also $1$. This proves the closed case, with the exact groups and the necessary conjugator adjustment specified.

An isomorphism of these two completions means exactly that finite-index congruence subgroups are cofinal among all finite-index subgroups. Restricting the same quotient-complex diagram to a cofinal index system gives the asserted isomorphism of profinite complexes. D.11 then gives the stated connectedness and fundamental groups. At complexity one the unmodified curve complex is disconnected; D.12 and D.21 settle complexity two. This argument asserts no unproved higher profinite connectivity.
\end{proof}

\begin{result}{D.20}{Refuted}{Two literal centralizer assertions in the retrieved Boggi source}
Lemma~6.7 of the retrieved 12 July 2021 version of \emph{Congruence topologies on the mapping class group} asserts that the hyperelliptic involution is self-centralizing in every finite covering-transformation semidirect product specified there. The characteristic mod-two homology cover refutes that assertion. Corollary~4.3 of the same version, when stated for arbitrary unequal multitwist exponents as printed, also cannot identify the centralizer with the full stabilizer of the underlying curve set. Neither assertion is used in D.19.
\end{result}
\begin{proof}
For a closed hyperelliptic surface of genus $g\geq2$, choose a Weierstrass basepoint and set
\[
 N=\ker\bigl(\pi_1(S_g)\longrightarrow H_1(S_g;\mathbb F_2)\bigr).
\]
This subgroup is proper, normal, finite index and characteristic, so it satisfies every subgroup hypothesis in the printed Lemma~6.7 \cite[pp.~26--27]{D_Boggi2020}. The hyperelliptic involution acts as $-I$ on integral homology and hence trivially on the quotient
$G_N=(C_2)^{2g}$. Thus the specified deck group is
\[
 G=G_N\rtimes\langle\iota\rangle
   =(C_2)^{2g+1},\qquad C_G(\iota)=G.
\]
Its centralizer has order $2^{2g+1}$, whereas $\langle\iota\rangle$ has order two. The unramified surface cover has degree $2^{2g}$ and genus $1+2^{2g}(g-1)$ by Euler characteristic. The saved exact computation gives degree $16$, genus $17$, and centralizer order $32$ at $g=2$.

For the second assertion choose disjoint nonseparating curves $a_1,a_2$ on the two handles of a closed genus-two surface, and an orientation-preserving mapping class $s$ exchanging the handles. Then $s$ belongs to the setwise stabilizer of $\{a_1,a_2\}$, which is the stabilizer explicitly defined in Theorem~3.6 of that source. But
\[
 s(T_{a_1}T_{a_2}^{\,2})s^{-1}
     =T_{a_1}^{\,2}T_{a_2}
     \ne T_{a_1}T_{a_2}^{\,2}.
\]
The inequality is visible already on $H_1(S_2;\mathbb F_3)$: in the symplectic basis $(a_1,b_1,a_2,b_2)$ the two matrices have upper-right entries $(1,2)$ and $(2,1)$ in their two diagonal $2$-by-$2$ blocks. Thus it persists in the procongruence group, without any appeal to a centralizer theorem. Both examples were executed in \path{code/d_hyperelliptic_counterexample.py}. The first refutes a universal finite-quotient statement; it does not refute the separate profinite self-normalizer assertion, which requires a cofinal argument. The second shows exactly why a weight-preserving stabilizer must appear when unequal exponents are used.
\end{proof}

\begin{result}{D.21}{Proved}{The second complexity-two profinite obstruction}
The full finite-level completion of the unmodified curve complex $\mathcal C_{1,2}$ has nonzero continuous first cohomology with $\mathbb F_2$ coefficients, and nontrivial profinite fundamental group. By D.19 the same holds for its procongruence completion.
\end{result}
\begin{proof}
Use the torus involution exchanging the two marked points. Its quotient is a sphere with four branch points and a fifth distinguished point, the image of the two marks. Lifting a sphere curve gives an isomorphism
\[
 P^*:\mathcal C_{0,5}\xrightarrow{\sim}\mathcal C_{1,2}.
\]
The precise map and invariant-representative proof are Luo, Lemma~2.1(a), opened on pages~3--4 \cite{D_Luo2000}. A sphere curve separating two branch points lifts to two parallel nonseparating curves and hence one isotopy class. A curve separating the distinguished point and one branch point lifts to one separating curve bounding the disk containing both torus marks. Hyperelliptic invariance supplies an invariant representative of every torus curve; taking invariant disjoint representatives proves the converse disjointness implication. These operations give the displayed inverse correspondence.

The construction is equivariant for the lift of the pure sphere mapping class group. Algebraically, B.9 identifies
\[
 \Gamma_{1,2}=B_4/\langle\Delta_4^2\rangle,\qquad
 \Gamma_{0,5}=P_4/\langle\Delta_4^2\rangle,
\]
and the latter subgroup has index $24$, by the permutation homomorphism to $\mathfrak S_4$. The subgroup is exactly the one retaining all four branch labels in the preceding cover construction. Finite-index subgroups of this subgroup are cofinal in the finite-index system for $\Gamma_{1,2}$; normality can be imposed by taking the finite intersection of conjugates. The equivariant simplicial isomorphism therefore identifies the two full finite-level curve complexes. Transporting the explicit continuous Petersen-edge class and the closed pentagon of D.12 gives a nonzero continuous $\mathbb F_2$ class on the genus-one complex. Its profinite fundamental group is consequently nontrivial. D.19 supplies the final procongruence comparison.
\end{proof}
